\section{The actual analytic functions}\label{sec:analytic}
All integrals in this section use the probability average
\[
 \langle h\rangle=\frac1{2\pi}\int_0^{2\pi}h(t)\,dt.
\]
In particular this is not an integral over a region selected by a pattern.
For each of the 43 rational coefficient records of Kapila's Appendix~A,
put
\[
 r_v(x)=\sum_{j=1}^8r_{v,j}\cos(jx),\quad
 s_v(x)=\sum_{j=1}^8s_{v,j}\cos(jx),\quad
 F_v=\exp((\ln2)r_v),\quad G_v=\exp((\ln2)s_v).
\]
Suppress $v$ temporarily. The exponent $\kappa=\kappa_v$ is the exact
source rational; $1\le\kappa\le2$. Define
\begin{align*}
 Z_f&=\langle F\rangle,& Z_g&=\langle G\rangle,&
 Z_{f,\kappa}&=\langle F^\kappa\rangle,&
 Z_{g,\kappa}&=\langle G^\kappa\rangle,\\
 f&=F/Z_f,&g&=G/Z_g,& f_*&=F^\kappa/Z_{f,\kappa},&
 g_*&=G^\kappa/Z_{g,\kappa}.
\end{align*}
\begin{align*}
 \mathcal P(x)&=\langle F(x+t)G(t)\rangle_t,\\
 \mathcal Y(x)&=\langle F(x+t)G(t)(r(x+t)+s(t))\rangle_t,\\
 Q(x)&=\mathcal P(x)/(Z_fZ_g),\\
 K_x(t)&=f(x+t)g(t)/Q(x).
\end{align*}
The intended scores, following the definitions rather than their names,
are
\begin{align*}
 D(x)&=\langle K_x\log_2K_x\rangle,&
 A_s(x)&=\langle K_x(t)\log_2f_*(x+t)\rangle_t,\\
 E_s(x)&=\langle K_x\log_2g_*\rangle,&
 u(x)&=A_s(x)+E_s(x)-D(x),\quad
 w(x)=2D(x)-A_s(x)-E_s(x).
\end{align*}
The subscripts distinguish scores $A_s,E_s$ from a budget ratio or an
exponent. All logarithms in scores are real base-two logarithms.

\begin{lemma}[Concrete existence and score bound]\label{lem:exist}
For these definitions all denominators are strictly positive, the four
density normalizations and $K_x$ have average one, and all real scores
are continuous and bounded. Every affine-certificate moment integrand
is continuous and integrable. If each coefficient has absolute value
at most two, then $|w(x)|\le256$ for every $x$.
\end{lemma}
\begin{proof}
The coefficient assumption gives $|r|,|s|\le16$. Hence
$2^{-16}\le F,G,Z_f,Z_g\le2^{16}$ and
$2^{-32}\le\mathcal P\le2^{32}$. The tilted quantities satisfy the
corresponding bounds with $16\kappa$ in place of 16.
All assertions of positivity follow before taking any logarithm or
quotient. Periodicity and substitution show that the averages of the
normalized functions are one. The average of $K_x$ is one by the
definition of $Q$. Integration of continuous functions depending
continuously on $x$ on a compact rectangle preserves continuity;
the strictly positive denominator bounds preserve it through division
and logarithms. This proves integrability as well as continuity, and
prevents a totalized logarithm or integral from masking a missing hypothesis.

Put
\[
 L=\log_2\mathcal P-\log_2(Z_fZ_g),\quad
 B=\mathcal Y/\mathcal P-\log_2(Z_fZ_g),\quad
 H=\log_2(Z_{f,\kappa}Z_{g,\kappa})-\kappa\log_2(Z_fZ_g).
\]
Here $B$ is the $K_x$-average of $r(x+t)+s(t)-\log_2(Z_fZ_g)$.
Substitution gives $D=B-L$ and $A_s+E_s=\kappa B-H$, so
\begin{equation}\label{eq:score-identities}
 u=L+(\kappa-1)B-H,\qquad w=(2-\kappa)B-2L+H.
\end{equation}
The quotient $\mathcal Y/\mathcal P$ is an average of $r+s$ with
positive weights summing to one, hence has absolute value at most 32.
Also $|\log_2(Z_fZ_g)|\le32$ and $|\log_2\mathcal P|\le32$.
Thus $|B|,|L|\le64$. The tilted normalizer bounds give $|H|\le64\kappa$.
Since $0\le2-\kappa\le1$,
\[
 |w|\le64(2-\kappa)+128+64\kappa=256.
\]
This derives the score bound through the actual definitions, rather
than stopping at a bound on trigonometric coefficients.
\end{proof}

\begin{lemma}[The required entropy order]\label{lem:entropy-order}
The actual scores satisfy $D\ge0$, $D-A_s\ge0$, $D-E_s\ge0$,
and hence $w\ge0$. The same is true for their exact finite-circle
versions.
\end{lemma}
\begin{proof}
If $p,q$ are strictly positive normalized densities, the elementary
inequality $\ln t\le t-1$ gives
\[
 \langle p\ln(q/p)\rangle\le\langle q-p\rangle=0.
\]
Apply this with $p=K_x$ and $q=1,f_*(x+\cdot),g_*$ respectively,
and divide by $\ln2>0$. Each $q$ is normalized; for the shifted
one this follows from periodic translation invariance. The same
argument uses only normalized sums on the finite circle, where grid
translations permute the summands. Finally
$w=(D-A_s)+(D-E_s)\ge0$.
\end{proof}

\paragraph{Independent normalization check.}
Eliminating all untilted normalizers from~\eqref{eq:score-identities}
gives the exact equivalent formulas
\begin{align}
 u&=\log_2\mathcal P+(\kappa-1)\mathcal Y/\mathcal P
       -\log_2(Z_{f,\kappa}Z_{g,\kappa}),\label{eq:cancel-u}\\
 w&=(2-\kappa)\mathcal Y/\mathcal P-2\log_2\mathcal P
       +\log_2(Z_{f,\kappa}Z_{g,\kappa}).\label{eq:cancel-w}
\end{align}
The tilted normalizers do not cancel. These identities are useful for
an independent implementation check, not alternative definitions that
avoid the density correspondence.

\section{A validated integration contract}\label{sec:quadrature}
This section proves a sufficient error bound. It does not assert that
any numerical inequality has passed until the specified enclosures
have actually been produced and checked. All source decimals are
exact rationals, including coefficients and multipliers.

\begin{lemma}[Periodic trapezoid bound]\label{lem:trap}
Suppose a $2\pi$-periodic function $h$ is holomorphic on a neighborhood
of the closed strip $|\operatorname{Im}z|\le\sigma$, where $\sigma>0$,
and $|h|\le B_0$ there. For $N\ge1$,
\[
 \left|\frac1N\sum_{j=0}^{N-1}h(2\pi j/N)-\langle h\rangle\right|
 \le \mathcal E_N(B_0,\sigma):=\frac{2B_0}{e^{N\sigma}-1}.
\]
\end{lemma}
\begin{proof}
For a nonzero Fourier index $k$, shift the integral for its coefficient
to $\operatorname{Im}z=-\sigma\operatorname{sgn}k$. Periodicity cancels
the vertical contour sides and gives
$|\widehat h(k)|\le B_0e^{-|k|\sigma}$. Thus the Fourier series is
absolutely convergent on the real axis. Averaging its terms at the
$N$ nodes kills every index not divisible by $N$. The error is at most
$2B_0\sum_{j\ge1}e^{-jN\sigma}$, the claimed geometric sum.
\end{proof}

\paragraph{Inner integrals and arithmetic.}
Define $R(t)=\sum_j|r_j|\cosh(jt)$ and
$S(t)=\sum_j|s_j|\cosh(jt)$. For an inner strip width $\sigma>0$,
$|r(z)|\le R(\sigma)$ and $|s(z)|\le S(\sigma)$.
Lemma~\ref{lem:trap} therefore applies with the following strip bounds:
\begin{center}\small
\begin{tabular}{ll}\toprule
Inner quantity & Valid $B_0$\\\midrule
$Z_f,Z_g$ & $2^{R(\sigma)},\ 2^{S(\sigma)}$\\
$Z_{f,\kappa},Z_{g,\kappa}$ & $2^{\kappa R(\sigma)},\ 2^{\kappa S(\sigma)}$\\
$\mathcal P(x_j)$ & $2^{R(\sigma)+S(\sigma)}$\\
$\mathcal Y(x_j)$ & $(R(\sigma)+S(\sigma))2^{R(\sigma)+S(\sigma)}$\\
$\langle2^{2r_\beta}\rangle$ &
 $2^{2\sum_j|r_j|\cosh(j\beta)\cosh(j\sigma)}$\\
$\langle2^{2s_\beta}\rangle$ &
 $2^{2\sum_j|s_j|\cosh(j\beta)\cosh(j\sigma)}$\\\bottomrule
\end{tabular}
\end{center}
Here $r_\beta(x)=\sum_jr_j\cosh(j\beta)\cos(jx)$; $s_\beta$ is
defined similarly. The node $x_j$ is real, so the same convolution bounds
hold uniformly for all outer nodes.
Outward interval arithmetic must enclose the exact node sums, including
the trigonometric arguments, exponentials, logarithms and all normalizers.
Enlarge each inner interval by its own $\mathcal E_N$ before performing
division or taking logarithms. Every denominator/logarithm interval used
must have a verified strictly positive lower endpoint. A larger precision
or agreement of two meshes is not a substitute for these error bounds.

\paragraph{The finite DFT acceleration is an exact sum identity.}
For node arrays $a_j=F(2\pi j/N)$ and $b_j=G(2\pi j/N)$ define
$\widehat a_\ell=\sum_j a_j e^{-2\pi i\ell j/N}$ and use the inverse
transform with prefactor $1/N$. Orthogonality of the finite geometric
sum gives, with all indices read modulo $N$,
\[
 \operatorname{DFT}^{-1}
       (\widehat a_\ell\widehat b_{-\ell})_j
 =\sum_t a_{j+t}b_t.
\]
Indeed, expansion of the left side retains exactly the terms with
$-j_a+j_b+j=0\pmod N$. Thus the raw convolution node average is this
inverse transform divided by $N$ once more. For $\mathcal Y$, add the
correlations of $(Fr,G)$ and $(F,Gs)$ and make the same division.
The identity applies to exact node values. Rigorous ball DFT arithmetic
must enclose that finite calculation; the analytic remainders above
still have to be added afterward. An ordinary floating FFT without
outward error control would not be a validated replacement.

\begin{lemma}[A genuine zero-free strip]\label{lem:strip}
Let $\beta>0$, and suppose rigorously computed bounds enclose
\[
 A_f=\frac{\langle2^{2r_\beta}\rangle}{Z_f^2},\quad
 A_g=\frac{\langle2^{2s_\beta}\rangle}{Z_g^2},\quad
 q_0=\sqrt{(A_f-1)(A_g-1)}<1.
\]
Then on $|\operatorname{Im}z|\le2\beta$,
\[
 |Q(z)-1|\le q_0,
 \qquad |L(z)|\le-\log_2(1-q_0),
\]
and $Q$ has a holomorphic logarithm real on the real axis. Moreover
\begin{equation}\label{eq:B-strip-bound}
 |B(z)|\le
 \frac{\sqrt{A_fA_g}(R(\beta)+S(\beta))}{1-q_0}
       +|\log_2(Z_fZ_g)|.
\end{equation}
\end{lemma}
\begin{proof}
The entire periodic functions $f,g$ are real and even on the real circle
because their defining trigonometric polynomials use only real cosines.
They have shifted means one by contour shifting. For a function real
on the real circle, its Fourier coefficients
satisfy $\widehat f(-k)=\overline{\widehat f(k)}$. Parseval therefore
writes the squared shifted $L^2$ norm as its constant term plus
nonnegative multiples of $\cosh(2ky)$, so it is nondecreasing in $|y|$.
This also bounds both positive and negative imaginary shifts.
At $y=\beta$ that squared norm is precisely $A_f$, since
$|F(t+i\beta)|^2=2^{2r_\beta(t)}$. In particular $A_f,A_g\ge1$.

For $z=x+iy$ with $|y|\le2\beta$, split the convolution contour:
\[
 Q(z)=\langle f(x+t+iy/2)g(t-iy/2)\rangle_t.
\]
Subtract the product of the shifted means and apply Cauchy--Schwarz
to $f-1,g-1$. Their squared norms are at most $A_f-1,A_g-1$;
Here the centered identity is
$\langle|f(\cdot+iy)-1|^2\rangle
=\langle|f(\cdot+iy)|^2\rangle-2\operatorname{Re}\langle f(\cdot+iy)\rangle+1
=\langle|f(\cdot+iy)|^2\rangle-1$, and similarly for $g$.
This proves the disk bound. The power series for $\operatorname{Log}(1+h)$,
with $h=Q-1$, supplies the holomorphic branch and gives
$|\operatorname{Log}Q|\le\sum_{j\ge1}q_0^j/j=-\ln(1-q_0)$.
Divide by $\ln2$ for the stated base-two bound. The strict inequality
$q_0<1$ also supplies a neighborhood of the closed strip.

Apply the same split contour to $\mathcal Y/(Z_fZ_g)$.
The absolute value of $r+s$ there is at most $R(\beta)+S(\beta)$;
Cauchy--Schwarz bounds the integral of $|fg|$ by $\sqrt{A_fA_g}$.
Since $|Q|\ge1-q_0$ and
$B=\mathcal Y/(Z_fZ_gQ)-\log_2(Z_fZ_g)$, this gives
\eqref{eq:B-strip-bound}. Real positivity alone would not justify
the complex logarithm used in this proof.
\end{proof}

\begin{proposition}[Per-certificate validated upper bound]\label{prop:validated}
For every filter occurring in a certificate or one of its comparison rows,
suppose Lemma~\ref{lem:strip}'s bounds are enclosed independently. Let
$L_v^*,B_v^*,H_v^*$ bound $|L_v|,|B_v|,|H_v|$ on the strip, and put
\[
 U_v^*=L_v^*+|\kappa_v-1|B_v^*+H_v^*,\qquad
 W_v^*=|2-\kappa_v|B_v^*+2L_v^*+H_v^*.
\]
For an old row vector $G_{v,i}$ form the following real-part upper bounds:
\begin{align*}
 \rho_0:&\quad (\Delta_b+U_v^*,\ \Gamma_b+W_v^*,\
                       (U_v^*+U_j^*)_{j\in C}),\\
 \rho_{1:8}:&\quad (A_bW_v^*+U_v^*,\ W_v^*-\Gamma_bL,\
                  \Gamma_bH+W_v^*,\ (U_v^*+U_j^*)_{j\in C}),\\
 \rho_9:&\quad (A_bW_v^*+U_v^*,\ W_v^*-15\Gamma_b,\
                       (U_v^*+U_j^*)_{j\in C}).
\end{align*}
Call the resulting vector $S_{v,i}$ and set
\[
 B_{v,i}^*=(2+2\cosh(2\beta))\,2^{z_{v,i}\cdot S_{v,i}}.
\]
If the interval $T_{v,i}$ encloses the average of the true outer node
values after all inner/arithmetic errors, then
\[
 I_{v,i}\le \sup T_{v,i}+\mathcal E_N(B_{v,i}^*,2\beta).
\]
Certification requires this upper bound to be strictly below an
independently enclosed lower bound for $2^{\eta_b-m_i}$.
\end{proposition}
\begin{proof}
Each affine row's real part is bounded by the displayed entry;
$A_b>0$ and all multipliers are nonnegative. Hence the real part of
the exponent is at most $z\cdot S$. The holomorphic extension uses
$2^z=\exp((\ln2)z)$ and the log branch already justified. Also
$|2+2\cos z|\le2+2\cosh(2\beta)$. Lemma~\ref{lem:trap}, followed
by the enclosing node interval, gives the result. Negative entries
such as $W_v^*-\Gamma_bL$ may be retained: they are valid upper bounds,
not assumed nonnegative rows at complex points.
\end{proof}

\paragraph{Exact target and pilot.}
The targets needed by the released theorem are $m_0=m_1=.00025$,
$m_2=\cdots=m_8=4$, and $m_9=17$, with $\eta_b=.51613$.
Kapila's common Proposition-B threshold alone is insufficient.
The bounded pilot consists of $(3,0),(3,1),(4,0)$ and at most one tail
certificate, using the actual source comparisons for each. A separate
record must classify every attempted pair as certified, unresolved or
refuted. A failed sufficient bound is unresolved, not a refutation;
a certified lower bound above threshold would be a refutation.
If bounds are rigorous and the measured cost permits the finite full
batch, the same contract can be applied to all 430 distinct pairs.

\section{Finite-circle convergence and effective evaluation}\label{sec:finite}
\begin{proposition}[All normalizers are part of the limit]\label{prop:finite}
Replace every inner average above by the normalized $P$-point average,
including the tilted normalizers. At the grid points the resulting
scores are precisely the finite-circle scores, and they converge
uniformly to the continuous scores. The corresponding outer finite
moments converge to their continuous moments. All convergence bounds
can be made effective for this fixed finite coefficient family.
\end{proposition}
\begin{proof}
For any continuously differentiable periodic real function $h$, comparison
on each interval of length $2\pi/P$ with its left endpoint gives
\[
 |T_Ph-\langle h\rangle|\le 2\pi\|h'\|_\infty/P.
\]
The same estimate holds uniformly in an additional real parameter
when its derivative bound is uniform. In this coefficient family,
$|r'|,|s'|\le2\sum_{j=1}^8j=72$. Thus valid bounds on the
integration-variable derivatives of $FG$ and $FG(r+s)$ are respectively
\[
 144(\ln2)2^{32},\qquad 144(32\ln2+1)2^{32}.
\]
For $F^\kappa,G^\kappa$ a bound is $144(\ln2)2^{32}$.
The raw normalizers and convolutions therefore differ from their
continuous versions by explicitly bounded constants divided by $P$.
Their positive lower bounds, for example $\mathcal P_P,\mathcal P\ge2^{-32}$
and $Z_{f,P},Z_f\ge2^{-16}$, hold on either circle without a limiting
argument. On these fixed positive compact intervals the quotient and
logarithm operations are Lipschitz; for the latter use
$|\log_2x-\log_2y|\le |x-y|/(\min(x,y)\ln2)$.
Substitution in~\eqref{eq:score-identities} gives computable constants
$C_u,C_w$ such that the uniform score errors are at most $C_u/P,C_w/P$.

For a fixed certificate, the exponent error is at most a computable
$C_E/P$. Its real exponents are uniformly bounded by another computable
constant $J$. The mean-value bound for $2^x$ on this interval, and
$0\le2+2\cos x\le4$, bound the change in the node integrand by a
constant divided by $P$. The continuous integrand is itself continuously
differentiable with a computable derivative bound, so applying the
Riemann estimate once more proves an explicit $C_I/P$ moment error.
All constants are finite expressions in the fixed rational coefficients,
$\pi,\ln2$ and elementary-function bounds, not input-dependent advice.
Consequently a certified continuous upper bound separated from a desired
finite threshold by $d>0$ gives the effective choice $P>C_I/d$.
There are finitely many modified certificates, so the maximum cutoff
works for all of them, with any strictly smaller finite margin.

For a prescribed fixed accuracy, choose a rational mesh from these
bounds and evaluate the continuous elementary functions by rational
Taylor series with explicit remainders. Bounded positive logarithm
arguments admit the convergent series for
$\ln y=2\sum_{j\ge0}((y-1)/(y+1))^{2j+1}/(2j+1)$.
Below the fixed cutoff compute finite averages directly; above it use
the continuous approximation and the uniform error bound. Approximating
the input residue angle uses polynomially many operations in $\log P$;
all meshes and accuracy constants are fixed. This gives an effective
uniform evaluation procedure, though not a claim that our Lean project
has already formalized this analytic argument.
\end{proof}

These arguments reconstruct the analytic steps in the pinned OpenAI
filtering/numerical sections and Kapila's finite-family adaptation
\cite{openai,kapila}. The proofs above supply their applicability to the
actual 43-filter definitions; the numerical moment thresholds remain
separate obligations until validated bounds are produced.
